\documentclass[11pt]{article}
\usepackage[T1]{fontenc}
\usepackage{amsmath,amssymb,amsthm}
\usepackage[margin=1in]{geometry}

\newtheorem{theorem}{Theorem}
\newtheorem{lemma}{Lemma}
\newtheorem{corollary}{Corollary}
\newtheorem{remark}{Remark}

\title{Laboratory V80: Hall Certificates and Support Branchwidth}
\author{NC0\_k-Avoid Laboratory}
\date{August 2026}

\begin{document}
\maketitle

Let $C:\{0,1\}^n\to\{0,1\}^m$ be a Boolean circuit. Let $M$ be its
output gates, let $X$ be the set of active input variables, and write
\[
 N(S)=\bigcup_{i\in S}\operatorname{supp}(i)
 \qquad(S\subseteq M).
\]
The support connectivity function is
\[
 \lambda_C(S)=|N(S)\cap N(M\setminus S)|.
\]

\begin{lemma}[Hall projection]
If $|N(S)|<|S|$, then the projected circuit on coordinates $S$ omits an
output pattern. Every full output word extending such a pattern is outside
$\operatorname{range}(C)$.
\end{lemma}

\begin{proof}
The projection depends on at most $|N(S)|$ Boolean variables, so its range has
size at most $2^{|N(S)|}<2^{|S|}$. A preimage of a full extension would give a
preimage of the omitted projected pattern.
\end{proof}

\begin{remark}[Algorithmic boundary]
If $|N(S)|=O(\log n)$, exhaustive evaluation constructs an omitted pattern in
polynomial time. For deficiency $d=|S|-|N(S)|\geq1$, a uniformly random pattern
on $S$ is omitted with probability at least $1-2^{-d}$; an NP oracle tests
range membership, giving a zero-error expected-polynomial randomized
procedure. This counting argument does not by itself give a deterministic
$\mathbf{FP}^{\mathbf{NP}}$ procedure or a polynomial canonical candidate
list.
\end{remark}

\begin{lemma}[Exact cut identity]
For every $S\subseteq M$,
\[
 \lambda_C(S)=|N(S)|+|N(M\setminus S)|-|X|.
\]
Consequently, if both $S$ and $M\setminus S$ are Hall-nondeficient, then
\[
 \lambda_C(S)\geq |M|-|X|.
\]
\end{lemma}

\begin{proof}
The first identity is inclusion--exclusion, since
$N(S)\cup N(M\setminus S)=X$. Under the two Hall inequalities, substitute
$|N(S)|\geq|S|$ and
$|N(M\setminus S)|\geq|M\setminus S|$.
\end{proof}

\begin{theorem}[Balanced Hall expansion forces support width]
Assume every $S\subseteq M$ with
$m/3\leq |S|\leq2m/3$ is Hall-nondeficient. Then
\[
 \operatorname{bw}(\lambda_C)\geq m-|X|.
\]
Equivalently, every branch decomposition of width below the stretch has a
balanced edge with a Hall-deficient side.
\end{theorem}

\begin{proof}
Every subcubic tree with $m$ labelled leaves has an edge for which both sides
contain between $m/3$ and $2m/3$ leaves. The two sides Hall-expand by
hypothesis, so the cut identity lower-bounds the width of that edge by
$m-|X|$.
\end{proof}

\begin{theorem}[Local Hall-expansion barrier]
For all sufficiently large $n$, with
$m=n+\lceil n^{2/3}\rceil$, there exists a rank-at-most-three support family
with no Hall-deficient gate set of size at most $n/(16e^2)$.
\end{theorem}

\begin{proof}
Choose three variables independently and uniformly for each gate, collapsing
repetitions. Fix $s\leq n/(16e^2)$ gates. A Hall violation is witnessed by
fewer than $s$ variables containing all $3s$ choices. Since $m\leq2n$ and
$s\leq n/2$, a union bound gives
\[
 \Pr[\text{a bad set of size }s]
 \leq {m\choose s}\,s{n\choose s}\left(\frac{s}{n}\right)^{3s}
 \leq s\left(\frac18\right)^s.
\]
Therefore the probability of any such bad set is at most
\[
 \sum_{s\geq1}s(1/8)^s=\frac8{49}<1.
\]
A family with the required property exists.
\end{proof}

\begin{remark}
The last theorem rules out a universal logarithmic-size Hall witness based only
on rank three and the target stretch. It does not prove high branchwidth for
every locally expanding family. The laboratory therefore targets deterministic
candidate lists from additional affine, sunflower, or overlap structure, with
explicit resistant families as the alternative output.
\end{remark}

\paragraph{Scientific boundary.}
No deterministic $\mathbf{FP}^{\mathbf{NP}}$ algorithm for unrestricted
$\mathrm{NC}^0_3$-Avoid, circuit lower bound, bounded-arithmetic formalization,
novelty claim, peer review, or P-versus-NP resolution is asserted.

\paragraph{References.}
K.~Gajulapalli, A.~Golovnev, S.~Nagargoje, and S.~Saraogi,
\emph{Range Avoidance for Constant-Depth Circuits: Hardness and Algorithms},
ECCC TR23-021, 2023. S.~Huang, X.~Li, and Y.~Zhong,
\emph{Range Avoidance and Remote Point: New Algorithms and Hardness},
ECCC TR25-049, 2025. H.~Ren, Y.~Wang, and Y.~Zhong,
\emph{Hardness of Range Avoidance and Proof Complexity Generators from
Demi-Bits}, ECCC TR25-191, 2025.

\end{document}
